% !TEX root = ../../main.tex
\chapter{Notation}
\label{ch:appendix:notation}

% BEGIN TUFTE PLACEHOLDER: supplied sample, lines 1153--1190.
\begingroup\sloppy
\label{template:sec:tl-messages}
The following is a list of all of the errors, warnings, and other messages generated by the \TL classes and a brief description of their meanings.
\index{error messages}\index{warning messages}\index{debug messages}

% Errors
\docmsg{Error: \doccmd{subparagraph} is undefined by this class.}{%
The \doccmd{subparagraph} command is not defined in the \TL document classes.
If you'd like to use the \doccmd{subparagraph} command, you'll need to redefine
it yourself.  See the ``Headings'' section on page~\pageref{template:sec:headings} for a
description of the heading styles availaboe in the \TL document classes.}

\docmsg{Error: \doccmd{subsubsection} is undefined by this class.}{%
The \doccmd{subsubsection} command is not defined in the \TL document classes.
If you'd like to use the \doccmd{subsubsection} command, you'll need to
redefine it yourself.  See the ``Headings'' section on
page~\pageref{template:sec:headings} for a description of the heading styles availaboe
in the \TL document classes.}

\docmsg{Error: You may only call \doccmd{morefloats} twice. See the\par\noindent\ \ \ \ \ \ \ \ Tufte-LaTeX documentation for other workarounds.}{%
\LaTeX{} allocates 18 slots for storing floats.  The first time
\doccmd{morefloats} is called, it allocates an additional 34 slots.  The second
time \doccmd{morefloats} is called, it allocates another 26 slots.

The \doccmd{morefloats} command may only be called two times.  Calling it a
third time will generate this error message.  See
page~\pageref{template:err:too-many-floats} for more information.}

% Warnings
\docmsg{Warning: Option `\docopt{class option}' is not supported -{}- ignoring option.}{%
This warning appears when you've tried to use \docopt{class option} with a \TL
document class, but \docopt{class option} isn't supported by the \TL document
class.  In this situation, \docopt{class option} is ignored.}

% Info / Debug messages
\docmsg{Info: The `\docclsopt{symmetric}' option implies `\docclsopt{twoside}'}{%
You specified the \docclsopt{symmetric} document class option.  This option automatically forces the \docclsopt{twoside} option as well.  See page~\pageref{template:clsopt:symmetric} for more information on the \docclsopt{symmetric} class option.}

\lipsum[1]
\par\endgroup
% END TUFTE PLACEHOLDER
